\documentclass[UTF8,a4paper,10pt]{ctexart}
\usepackage{ipara-handbook}
\usepackage{booktabs,tabularx,tikz}
\usetikzlibrary{arrows.meta,positioning}
\begin{document}
\section*{H-B02\quad 局部模型与区间定理：中值点、余项与误差控制}
\textbf{边界：}本桥只负责把局部近似升级为区间结论。Taylor 多项式与中值定理的完整机制仍由 Topic03、Topic04 持有。
\begin{tcolorbox}[title=母接口]
目标误差 $\longrightarrow$ 减去低阶模型 $\longrightarrow$ 构造零点函数
$\longrightarrow$ Rolle 链 $\longrightarrow$ 中值点余项/导数估计。
\end{tcolorbox}
\begin{center}
\begin{tikzpicture}[node distance=8mm, every node/.style={draw,rounded corners,align=center,minimum width=3.1cm,minimum height=8mm,font=\small}, >=Latex]
\node (a) {局部模型\\$P_n(x)$};
\node (b) [right=of a] {消去已知零点\\$\Phi=f-P_n-K\Pi$};
\node (c) [right=of b] {Rolle 链\\重复取导};
\node (d) [right=of c] {余项\\误差界};
\draw[->,thick] (a)--(b); \draw[->,thick] (b)--(c); \draw[->,thick] (c)--(d);
\end{tikzpicture}
\end{center}
\subsection*{1.\ Taylor 余项的接口}
若 $f\in C^{n+1}$，则在区间内存在 $\xi$ 使
\[
f(x)=\sum_{k=0}^{n}\frac{f^{(k)}(a)}{k!}(x-a)^k+
\frac{f^{(n+1)}(\xi)}{(n+1)!}(x-a)^{n+1}.
\]
因此若 $|f^{(n+1)}(t)|\le M$，立即得到
\[
|R_n(x)|\le \frac{M}{(n+1)!}|x-a|^{n+1}.
\]
这里的停止条件是：学生能说明“$\xi$ 只保证存在，$M$ 来自整个区间上界”，即可进入数值估计。
\subsection*{2.\ 多中值点：先制造插点，再分割区间}
中值定理在同一个区间反复使用，只能保证“各自存在”，不能保证得到的中值点不同。若目标需要 $\xi_1<\xi_2$ 等多个见证点，稳定做法是先反推一个插点 $c$ 应满足的函数值条件，用介值/零点定理证明 $c$ 存在，再在互不重叠的子区间上分别调用中值定理：
\[
[a,b]\longrightarrow [a,c]\cup[c,b]
\longrightarrow \xi_1\in(a,c),\quad \xi_2\in(c,b).
\]
多个中值点的“不同”来自\textbf{区间分离}，不是来自“同一条定理用了两次”。若要三个见证点，就先设计两个插点并形成三个子区间。

\subsection*{3.\ 中值点参数：存在性怎样升级为渐近位置？}
对每个充分小的 $h\ne0$，Lagrange 给出某个 $\theta_h\in(0,1)$ 使
\[
f(a+h)-f(a)=h f'(a+\theta_h h).
\]
定理本身不保证 $\theta_h$ 唯一，也不保证它有极限；要定位它，必须把\textbf{区间平均变化率}和\textbf{局部导数模型}放到同一阶比较。

若 $f''(a)\ne0$ 且二阶正则性足够，则
\[
\frac{f(a+h)-f(a)}h=f'(a)+\frac12f''(a)h+o(h),
\]
而
\[
f'(a+\theta_hh)=f'(a)+f''(a)\theta_hh+o(h),
\]
故
\[
\boxed{\theta_h\to\frac12}.
\]
更一般地，若第一个非零高阶导数是 $f^{(m)}(a)$（$m\ge2$），即
\[
f''(a)=\cdots=f^{(m-1)}(a)=0,\qquad f^{(m)}(a)\ne0,
\]
则比较 $h^{m-1}$ 主部得到
\[
\boxed{\theta_h\to m^{-1/(m-1)}}.
\]
若 $f$ 在附近仿射，中值等式对任意 $\theta_h\in(0,1)$ 都成立，说明没有非退化高阶信息时，中值点位置可能根本不被确定。

\subsection*{4.\ 插值余项：零点与重零点怎样编码约束？}
取互异节点 $x_0,\ldots,x_n$，令 $\Pi(x)=\prod_{i=0}^{n}(x-x_i)$。若目标点 $x$ 恰好等于某个节点，则 $f(x)-P(x)=0$，余项结论平凡成立，\textbf{不要再为了构造 $K$ 去除以 $\Pi(x)=0$}。以下只讨论目标点不是节点的情形。对
\[
\Phi(t)=f(t)-P(t)-K\Pi(t)
\]
选择 $K$ 使 $\Phi(x)=0$，并利用节点零点与 Rolle 定理逐层取导，得到 Lagrange 插值余项
\[
f(x)-P(x)=\frac{f^{(n+1)}(\xi)}{(n+1)!}\Pi(x).
\]
这不是一条孤立公式，而是
\[
\boxed{\text{插值约束}\to\text{制造零点}\to\text{反复 Rolle}\to\text{高阶导数见证}}.
\]

若节点还匹配导数，例如同时要求 $P(x_i)=f(x_i)$ 与 $P'(x_i)=f'(x_i)$，则 $f-P$ 在该节点具有重零点。Hermite 余项只是把“零点个数”升级为“按重数计数”；每增加一个独立导数约束，就为 Rolle 链多提供一层消失条件。完整高阶插值理论不进入主干，但这个\textbf{约束数 $\leftrightarrow$ 零点重数 $\leftrightarrow$ 余项阶数}接口必须保留。

\subsection*{5.\ 从点态余项到积分误差与导数估计}
若已有统一点态误差
\[
|f(x)-P(x)|\le E(x),
\]
则在合法区间上可直接累积为
\[
\left|\int_a^b(f-P)\right|\le\int_a^b E(x)\,dx.
\]
所以“先造余项，再积分”往往比直接攻击复杂积分更稳定。

端点零值还可以与 Rolle/Taylor 共同产生全区间界。例如若 $f(0)=f(1)=0$ 且 $|f''|\le M$，则对每个 $x\in[0,1]$ 可构造二次插值余项得到
\[
|f(x)|\le\frac M2x(1-x),
\qquad
\left|\int_0^1f(x)\,dx\right|\le\frac M{12}.
\]
这里 $x(1-x)$ 不是要背的模板，而是两个端点零约束生成的二次零点因子。

反向地，局部 Taylor 也能把函数值界与高阶导数界运输成中间导数界。若 $|f|\le M_0$、$|f''|\le M_2$，且 $[x-h,x+h]$ 都在合法区间内，则中央差分与 Taylor 给出一种典型估计
\[
|f'(x)|\le \frac{M_0}{h}+\frac{M_2h}{2}.
\]
选择合适的 $h$ 是在“函数值振幅”与“曲率误差”之间做平衡；完整 Landau--Kolmogorov 不等式只作 Extension。

\subsection*{6.\ 滑动区间 Taylor：无穷远处的中间导数怎样被夹住？}
若
\[
f(x)\to L,\qquad f''(x)\to0\qquad(x\to+\infty),
\]
固定一个 $h>0$，在 $[x,x+h]$ 上写 Taylor：
\[
f(x+h)=f(x)+hf'(x)+\frac{h^2}{2}f''(\xi_x).
\]
于是
\[
f'(x)=\frac{f(x+h)-f(x)}h-\frac h2f''(\xi_x)\to0.
\]
这条机制说明“低阶函数本身趋稳 + 更高阶变化率趋零”可以通过一个\textbf{随 $x$ 滑动但长度固定的区间}控制中间导数。高阶版本仍是同一思想：用足够阶 Taylor 把目标导数夹在端点差与更高阶余项之间。

\subsection*{7.\ 最小例子与调用}
对 $f(x)=e^x$ 在 $a=0$ 作一阶 Taylor，若 $|x|\le1$，存在 $\xi$ 介于 $0,x$ 之间使
\[
e^x=1+x+\frac{e^\xi}{2}x^2,\qquad
|e^x-1-x|\le\frac e2x^2.
\]
局部系数来自 Topic03，区间导数上界和中值点余项由本 Bridge 接通。题面要求误差界或“存在 $\xi$”时调用。

\subsection*{8.\ 合法性与反例边界}
\begin{itemize}
\item 中值点一般不能固定为端点、中心或某个“看起来合理”的点；
\item 多个中值点必须用区间分离保证不同，不能把“存在若干点”偷偷写成同一点或自动不同；
\item $\theta_h$ 的位置需要第一个非零高阶信息；仿射退化时它可以任意漂移；
\item 只有局部导数信息时，不能推出整个区间的统一误差界；导数上界的适用区间必须覆盖全部中值点；
\item 插值节点若重合，必须按重零点/导数条件重新计数，不能仍套普通 Lagrange 节点公式。
\end{itemize}
\begin{tcolorbox}[title=检查句]
“我用的是哪条中值定理？零点在哪里？导数上界覆盖哪个区间？余项中的 $\xi$ 是否仍只是存在性变量？”
\end{tcolorbox}
\end{document}
